\section{线性系统与矩阵指数}
\label{sec:15-Differential-Equations-and-Dynamical-Systems/01-ODE-Theory/02}

你手上有一个八维状态的线性仿真，把它积分到 \(t=200\)，输出全是 \(10^{30}\) 量级的数，最后变成 NaN。换一组初值重跑，曲线平稳地压向零。同事问：这是数值方法炸了，还是模型本身就该炸？

这个问题在跑仿真之前就能回答，而且只需要看系数矩阵的特征值。本节要建立的就是这条通道：把一维的 \(x'=kx\) 抬到 \(n\) 维，解写成矩阵指数，长期行为写在谱上。

\subsection{一维的答案能不能照抄到 \(n\) 维}

上一节末尾，谐振子被写成了矩阵形式的方程组。一般的 \(n\) 维线性系统写成

\[
x'=Ax,
\qquad
x(t)\in\R^n,\quad A\in\R^{n\times n},
\qquad
x(0)=x_0.
\]

一维情形 \(x'=kx\) 的解是 \(x(t)=e^{kt}x_0\)。把 \(k\) 换成矩阵 \(A\)，答案能不能照抄成

\[
x(t)=e^{tA}x_0\,?
\]

敢这样写，先得说清 \(e\) 的矩阵次幂是什么意思。

\subsection{矩阵指数就是那个幂级数}

数 \(e^z\) 的定义级数只用到乘法与加法：

\[
e^{z}=\sum_{k=0}^{\infty}\frac{z^k}{k!}.
\]

矩阵的幂与线性组合都有定义，同一个级数因此可以整段搬进矩阵世界：

\[
e^{tA}=\sum_{k=0}^{\infty}\frac{t^k}{k!}A^k.
\]

收敛性不成问题。由 \(\norm{A^k}\le\norm{A}^k\)，各项被收敛的数项级数 \(\sum_k(t\norm{A})^k/k!=e^{t\norm{A}}\) 控制，比较判别给出收敛，且所得级数可以对 \(t\) 逐项求导。这一步是\hyperref[sec:02-Calculus/05-Sequences-and-Series/05]{``幂级数是一族函数''}中理论的直接沿用，公式照抄，结论照用。

这个对象在本书已经露过面。\hyperref[sec:03-Linear-Algebra/05-Eigenvalues-and-Eigenvectors/04]{``矩阵函数：从多项式到指数''}用同一个幂级数把标量函数搬到矩阵上；\hyperref[sec:03-Linear-Algebra/05-Eigenvalues-and-Eigenvectors/02]{``对角化与动力系统：把耦合演化变成独立缩放''}则指出离散迭代 \(x_{k+1}=Ax_k\) 的解 \(x_k=A^kx_0\) 与连续方程解之间的平行：离散看矩阵的幂，连续看矩阵的指数。这里补上微分方程视角缺的那一步，逐项求导，

\[
\frac{d}{dt}e^{tA}
=\sum_{k=1}^{\infty}\frac{t^{k-1}}{(k-1)!}A^k
=Ae^{tA},
\]

于是 \(x(t)=e^{tA}x_0\) 满足 \(x'=Ax\)，而 \(t=0\) 时 \(e^{0}=I\) 给出正确初值。照抄的答案合法。

\subsection{可对角化时，解是特征模态的叠加}

写成 \(e^{tA}x_0\) 还谈不上``解出来了''的感觉——你没法从这个符号里读出仿真会不会炸。设 \(A\) 可对角化，\(A=S\Lambda S^{-1}\)，其中 \(\Lambda=\diag(\lambda_1,\ldots,\lambda_n)\)，\(S\) 的列是特征向量 \(v_i\)。由 \(A^k=S\Lambda^kS^{-1}\)，级数每一项都带着同一对 \(S,S^{-1}\)，提出来即得

\[
e^{tA}=Se^{t\Lambda}S^{-1},
\qquad
e^{t\Lambda}=\diag(e^{\lambda_1t},\ldots,e^{\lambda_nt}).
\]

把初值分解到特征基上，\(x_0=c_1v_1+\cdots+c_nv_n\)，解立刻变成

\[
x(t)=\sum_{i=1}^{n}c_i\,e^{\lambda_it}v_i.
\]

耦合的演化被拆成一组互不干扰的标量指数。每个特征方向是一个独立模态，第 \(i\) 个模态按 \(e^{\lambda_it}\) 增长或衰减。在可对角化的情形下，矩阵指数的全部内容就是这句话：特征值决定每个方向的指数率。开头那个仿真的谜底也在这里——某个 \(\lambda_i\) 的实部为正，而换初值时恰好把对应的 \(c_i\) 取成了零。

特征值可以是复数。写 \(\lambda=\alpha+i\beta\)，Euler 公式给出

\[
e^{\lambda t}=e^{\alpha t}(\cos\beta t+i\sin\beta t),
\]

实部 \(\alpha\) 控制振幅的指数增减，虚部 \(\beta\) 控制旋转的快慢。实矩阵的复特征值成共轭对出现，两个复模态合成实平面内的螺旋。上一节的谐振子是 \(\alpha=0\) 的特例：矩阵 \(\begin{pmatrix}0&1\\-\omega^2&0\end{pmatrix}\) 的特征值为 \(\pm i\omega\)，纯旋转，不增不减。

\subsection{判据是同一个，只是换了坐标}

于是长期行为的判据可以画在复平面上。

\begin{center}
\begin{tikzpicture}[>=stealth]
  \fill[gray!15] (-1.8,-1.45) rectangle (0,1.45);
  \draw[->] (-1.9,0) -- (1.9,0) node[right] {\(\operatorname{Re}\)};
  \draw[->] (0,-1.55) -- (0,1.6) node[above] {\(\operatorname{Im}\)};
  \fill (-1.3,0) circle (2pt);
  \fill (-0.7,0.8) circle (2pt);
  \fill (-0.7,-0.8) circle (2pt);
  \fill (0,1.1) circle (2pt);
  \fill (0,-1.1) circle (2pt);
  \fill (0.49,0.48) circle (2pt);
  \fill (0.49,-0.48) circle (2pt);
  \node at (-1.1,-1.15) {\small 衰减};
  \node at (1.15,-1.15) {\small 爆炸};
  \node[above right] at (0.05,1.15) {\small 振荡};
  \node at (0,-1.95) {\(x'=Ax\)：看 \(\operatorname{Re}\lambda\) 的符号};

  \draw[->] (2.7,0) -- (3.6,0);
  \node[above] at (3.15,0.05) {\small \(\lambda\mapsto e^{\lambda\Delta t}\)};

  \begin{scope}[shift={(5.6,0)}]
    \fill[gray!15] (0,0) circle (1);
    \draw (0,0) circle (1);
    \draw[->] (-1.9,0) -- (1.9,0) node[right] {\(\operatorname{Re}\)};
    \draw[->] (0,-1.55) -- (0,1.6) node[above] {\(\operatorname{Im}\)};
    \fill (0.27,0) circle (2pt);
    \fill (0.35,0.35) circle (2pt);
    \fill (0.35,-0.35) circle (2pt);
    \fill (0.45,0.89) circle (2pt);
    \fill (0.45,-0.89) circle (2pt);
    \fill (1.45,0.75) circle (2pt);
    \fill (1.45,-0.75) circle (2pt);
    \node at (-1.05,-1.15) {\small 衰减};
    \node at (1.4,-1.15) {\small 爆炸};
    \node at (0,-1.95) {\(x_{k+1}=Bx_k\)：看 \(\abs{\mu}\) 与 \(1\)};
  \end{scope}
\end{tikzpicture}
\end{center}

\emph{连续时间的稳定区是左半平面，离散时间的稳定区是单位圆盘；指数映射把前者送入后者，但会沿虚方向周期性卷绕。}

图右侧的判据来自离散迭代 \(x_k=B^kx_0\)：模长小于 \(1\) 的特征值让对应模态逐步消失，模长大于 \(1\) 的把它撑大。图左侧的判据来自 \(e^{\lambda t}\) 的模 \(e^{\alpha t}\)。两者由 \(\abs{e^{\lambda\Delta t}}=e^{\alpha\Delta t}\) 连起来：\(\alpha<0\) 当且仅当 \(\abs{e^{\lambda\Delta t}}<1\)。因此指数映射让两套稳定判据精确对应；但这不是一一换坐标，因为相差 \(2\pi i/\Delta t\) 的连续特征值会映到同一个离散特征值，而且有限的 \(\lambda\) 不会映到圆心 \(0\)。这也解释了为什么循环网络里``隐状态权重的谱半径超过 \(1\) 就会发散''和线性系统里``特征值跑到右半平面就会爆炸''听上去是两句话，用起来是一件事。

二维的相图分类现在只是这张图的读法。两个负实特征值时所有方向一起衰减，轨迹几乎沿直线沉入原点，称稳定结点；两个正实特征值把轨迹一律推开，称不稳定结点；一正一负则一个方向吸入、一个方向抛出，得到鞍点。复共轭对带来旋转：实部为负是稳定螺旋，实部为正是不稳定螺旋，落在虚轴上则振幅不变，称中心。这些名字此刻还只是特征值组合的几何翻译，下一单元\hyperref[sec:15-Differential-Equations-and-Dynamical-Systems/02-Stability-and-Dynamical-Systems/01]{``相平面：把解画成几何''}会把它们真的画出来，并追问非线性系统在多大程度上继承这张表。

顺带说一句边界。若 \(A\) 不可对角化，模态之间不再完全解耦，Jordan 块会在解里带出 \(t^{m}e^{\lambda t}\) 这样的多项式因子；但只要 \(\operatorname{Re}\lambda<0\)，指数仍然压得住多项式，衰减的结论不变。落在虚轴上的重特征值才是真正微妙的情形，那里多项式因子没有指数压制，系统会缓慢发散。

\subsection{有外界输入时，解是对过去的指数加权记忆}

若系统还受外界驱动，

\[
x'=Ax+b(t),
\]

解可以用同一个矩阵指数写成

\[
x(t)=e^{tA}x_0+\int_0^t e^{(t-s)A}b(s)\,ds.
\]

这里只取它的读法而不展开推导。第一项是初值的自由演化；积分项把每个时刻 \(s\) 的输入 \(b(s)\) 用剩下的时间 \(t-s\) 按同一矩阵指数传播到当下，再全部叠加。线性系统对外界的记忆，就是对过去输入的指数加权累计。

冬天供暖时每天都能碰上这种记忆。设 \(x(t)\) 记录两间相邻房间相对室外的温度，矩阵 \(A\) 描述房间向室外散热以及隔墙两侧的热交换，\(b(t)\) 记录暖气在各时刻的供热功率。两个房间此刻把暖气开到同一档，室温也可能差很远，因为开机前的温度和此前几小时的供暖过程都还在起作用：前者留在第一项，后者累积在积分里。矩阵指数算的正是每一段旧热量经过一段时间后还剩多少影响。

同一结构也出现在更专业的场景里。线性两室药代模型中，房间换成血浆与组织，散热换成药物清除，暖气换成输注泵的给药历史。公式一字不改，改的只是状态与参数的现实含义。若清除率或室间交换被错设，积分算得再准，也只是对错误生理模型的精确记忆。

\subsection{本节的落点}

线性系统被矩阵指数完全解决，而解的每个定性特征——增长、衰减、旋转、哪个方向占主导——都能从谱上直接读出。这张``由谱读行为''的样板值钱在后面：下一单元处理非线性系统的策略，正是在每个平衡点附近换成线性系统，再用本节的语言读出局部行为。本节还附赠一个全局结论——\(f(x)=Ax\) 满足最整齐的光滑条件，解对所有时间存在且唯一。下一节把``线性''换成一般函数，这两条都会失守。
